\section{Introduction}
\label{sec:introduction}

In a compact Calabi--Yau threefold, distinct singular regions
can couple to the same vector fields. This changes which
combinations of their local Higgs directions are available.
At a conifold, the relevant degrees of freedom are charged
wrapped-M2 hypermultiplets
\cite{Strominger:1995Conifold,Greene:1995hu,Katz:1996}.
At a collapsing del Pezzo surface they instead form an
interacting five-dimensional fixed point
\cite{Seiberg:1996,Morrison:1996pp,Douglas:1996,Intriligator:1997}.
For a degeneration containing both kinds of sector, the
compactification gauges flavor Cartans of the interacting
theory together with the charges of the nodal hypermultiplets.
The resulting constraints need not separate into independent
local transitions.

We study this mechanism and the vector-multiplet couplings
across one explicit compact transition. Its singular member
$X_9$ has two ordinary double points and a cone over
$\operatorname{Bl}_2\mathbb P^2$, which engineers the
rank-one $E_2$ theory. In a crepant resolution $\wideX_9$,
let $C_1,C_2$ be the two exceptional nodal curves and let
$R$ be the image of the $SU(2)$ flavor root of the del Pezzo
surface. The basic relation is integral:
\begin{equation}
 R=C_1+C_2\quad\hbox{in }H_2(\wideX_9;\ZZ).
\label{eq:intro-root-relation}
\end{equation}
The two independent compact gaugings therefore leave a single
joint Higgs direction. No local sector can move alone.
We use \emph{cooperative Higgsing} for this situation: a
global deformation activates several local sectors, but no
nonzero deformation supported on any one of them is allowed.
The explicit smoothing exists to all orders and changes
the Hodge numbers as
\begin{equation}
 (h^{1,1},h^{2,1})=(6,122)\longrightarrow(3,123).
\label{eq:intro-hodge-change}
\end{equation}

Our first result identifies the compact charge constraint and
its physical scope. Wrapped-M2 charges and flavor-root
pairings define a common linear map $Q$. On a chosen local
Higgs component, the image of the rigid triplet-moment-map
quotient in invariant deformation coordinates is $\ker Q$.
For the admissible toric hypersurfaces specified below, this
is also the image of the global first-order deformation map.
The exact mixed smoothing theorem further makes this a necessary
and sufficient existence test after the local discriminants are
excluded. Its selected-base integrability statement and the
explicit construction of the smoothing are established in
\cite{Wuebben:2026smoothing}; we recall their hypotheses and the
local identifications. Thus existence of the $X_9$ transition is
both a consequence of the general criterion and an explicit
construction used to determine its smooth phase. The invariant image is not the full
quaternionic-K\"ahler hypermultiplet manifold of the compact
supergravity theory. In particular, a rigid quotient does not
by itself determine its finite-Planck metric
\cite{Mohaupt:2004de}.

For $X_9$, the rank-two constraint accounts for two Higgsed
vectors; a third vector is lost when the $E_2$ Coulomb
direction disappears. The surviving charge lattice is computed
integrally, not only over $\QQ$. Its cubic and curvature
couplings agree with an independently constructed smooth
threefold. Locally a contractible $\mathbb F_1$ is replaced
by a rigid $\mathbb P^2$, giving the $E_0$ sector on the
smooth side. This calculation also distinguishes the primitive
unbroken torus from the cyclic group appearing in the formal
invariant ring.

Our second result concerns decoupling gravity. We classify the
nef directions in which the total volume diverges while the
nodal masses and the entire del Pezzo K\"ahler class remain
bounded. Both vectors imposing the joint constraint become
non-dynamical in every such limit. One boundary direction
can retain a different nodal gauging; no direction retains the
$E_2$ flavor-root gauging. Thus this particular coupling of
the local sectors requires finite Planck mass, under the stated
local-scale hypotheses. The conclusion follows from the
compact intersection form and includes a calculation with
the full inverse kinetic matrix.

The third and most detailed result determines the genus-zero
quantum vector geometry across the transition. On an auxiliary
circle, the quantum transition is a locus with four ordinary
conifold states spanning a primitive rank-three electric
lattice. An explicit residue-preserving Laurent mutation
relates its mirror to that of the smooth phase. We continue
the periods to this locus, prove descent of the seven regular
D0, D2 and D4 periods, and identify the eighth period in an
integral topological D-brane frame. This fixes the reduced
prepotential, including its quadratic ambiguity and constant
term, and hence the complete genus-zero special K\"ahler
geometry on the marked branch. The calculation is not merely
a comparison of classical cubics or a finite list of
Gromov--Witten invariants.

The finite-radius description has an important limitation.
The four conifold states reproduce the change
\eqref{eq:intro-hodge-change}, but their four-dimensional
effective theory does not remain uniform as the auxiliary
circle decompactifies. A correlated path reaches the
five-dimensional transition with finite bulk kinetic terms;
the Kaluza--Klein gap nevertheless closes. This separates
the exact period comparison from an unjustified claim that
four hypermultiplets give a complete five-dimensional
description of the strongly coupled transition.

The fibration structures provide independent physical checks.
One gauging vector is the K3-fiber class, identified with
the heterotic dilaton multiplet. A flat elliptic phase has a
six-dimensional $U(2)$ theory with anomaly-consistent matter
and a geometrically determined height pairing. The relation
\eqref{eq:intro-root-relation} identifies the images of roots
in two E-string lattices. An exact volume bound excludes the
six-dimensional F-theory limit on the transition face. The
heterotic instanton and five-brane counts and an explicit
candidate K3 are determined, but the global bundle construction
and its massless $U(1)$ remain open.

\paragraph{Relation to earlier work.}
The charged-hypermultiplet description of conifold transitions
and their supergravity couplings is standard
\cite{Greene:1995hu,Mohaupt:2004de}. Compact del Pezzo
transitions, including examples involving additional conifold
singularities, were studied in \cite{Malyshev:2007zz}.
The local Higgs rings used here are those of
\cite{Cremonesi:2015lsa}; the deformation rings follow from
Altmann's theory \cite{Altmann:1997}.
Conifold transitions relate variations of Hodge structure and
quantum cohomology in a general framework
\cite{Lee:2018Conifold}, and the toric Frobenius method is
well established \cite{Hosono:1995Mirror,Demirtas:2024Mirror}.
Our contribution is the explicit combination of a mixed
compact charge constraint, its local-scale decoupling
analysis, and an integral analytic period comparison for the
same transition. We do not propose a new general mirror
algorithm or a general theorem that all such mixed
transitions are smoothable.

\paragraph{Organization.}
Section~\ref{sec:compact-gauging} states the compact constraint
and its supergravity qualification.
Section~\ref{subsec:x9} gives the transition and integral
classical couplings, and Section~\ref{sec:decoupling} studies
gravity-decoupling limits. Section~\ref{sec:quantum} develops
the quantum transition and period comparison.
Section~\ref{sec:fibrations} gives the fibration interpretation.
The appendices contain the deformation, lattice, analytic
continuation and fibration calculations. Magnetic sheaves and
their stability walls are treated separately in
\cite{Wuebben:2026magnetic}; none of the period proofs below
depends on that investigation.
